\documentclass[../principal.tex]{subfiles}
\graphicspath{{\subfix{../imagenes/}}}

\begin{document}

\chapter{Lienzo en p5}

En p5.js dibujaremos en un lienzo, una superficie rectangular en 2 dimensiones. El lienzo en p5.js recibe el nombre de canvas.

\section{Contexto}

El lienzo es una abstracción del elemento canvas introducido en HTML 5.

En Processing existe una ventana donde renderizamos las gráficas que programamos, y donde también podemos programar interactividad con interfaces humano-computador como un teclado o un ratón. En p5.js, este lugar recibe el nombre de lienzo o canvas en inglés, ya que internamente está implementado con un elemento canvas de HTML 5.

\section{Sintaxis en p5.js}

En p5.js usamos la función createCanvas() para crear un lienzo, y que tiene los parámetros ancho (en eje x / horizontal), y altura (en eje y / vertical), ambas medidas en pixeles.

\begin{lstlisting}[language=JavaScript]
    createCanvas(ancho, altura);
\end{lstlisting}

\section{Referencias}

\begin{enumerate}
    \item  \url{https://en.wikipedia.org/wiki/Canvas_element}
\end{enumerate}

\end{document}


